\documentclass[../../../thesis.tex]{subfiles}
\begin{document}
\subsection{Formát dokumentu}
Rozmery stránky, typy písma, veľkosti, riadkovanie,
medzery medzi odsekmi, formát nadpisov, obrázkov, tabuliek,
rovníc a~ďalšie vizuálne parametre záverečnej práce
rešpektujú do maximálnej miery normu STN 01 6910: 2023
Pravidlá písania a úpravy písomností~\cite{stn016910}.

\subsubsection*{Rozmery strany}
Veľkosť bežnej textovej strany záverečnej práce je A4,
t.~j. 21\,cm $\times$ 29,7\,cm.
Pravý a~ľavý okraj majú šírku 2,75\,cm,
horný a~dolný okraj majú výšku 3\,cm.
Päta stránky, v~ktorej sa nachádza číslo strany,
je od spodnej hrany stránky vzdialená o~1,25\,cm.
Šírka textu je 15,5\,cm, jeho výška 23,7\,cm.
Horný a~dolný okraj obálky sú z~estetických
dôvodov zmenšené na 2\,cm.

\subsubsection*{Písmo a riadkovanie}
Základný font šablóny je normálny rez tzv. antikvového písma
s~veľkosťou 12\,pt.
V~tejto šablóne je to Computer Modern.
Vhodné sú aj iné fonty s pätkami ako Times, Georgia, Palatino a~podobne.
Na obálke a~titulnom liste používame bezpätkový (grotesk) font Latin Modern.
Jednotlivé typy odsekov (nadpisy, poznámky a~pod.)
majú jednotný typ písma,
odlišnosti vyjadrujeme rezom (polotučné písmo, kurzíva)
alebo veľkosťou.

Parameter \verb|\linespread| má hodnotu 1,25, t.~j.
vzdialenosť riadkov textu vo veľkosti 12\,pt je 15,6\,pt.

\subsubsection*{Nadpisy}
Šablóna záverečnej práce FEIstyle je založená na
štandardnej šablóne \LaTeX-u article.
Nadpis najvyššej úrovne je \verb|\section| zodpovedajúci kapitole.
Podkapitoly sú \verb|\subsection| a~\verb|\subsubsection|.
Číslovanie kapitol a podkapitol je viacúrovňové typu X.Y.Z, kde X je číslo kapitoly, Y je číslo podkapitoly a Z je číslo časti podkapitoly.
Číslovanie vyšších úrovní nie je definované.
Tvar a forma nadpisov zodpovedá norme STN ISO 2145: 1978 Dokumentácia.
Číslovanie oddielov a~pododdielov písaných dokumentov~\cite{iso2145}.

Nová kapitola začína vždy na novej strane.
Príkaz \verb|\section| spôsobí okrem sadzby čísla a názvu kapitoly aj ukončenie predošlej kapitoly, vysádzanie všetkých plávajúcich objektov (obrázky, tabuľky, výpisy kódu), ktoré sa nepodarilo umiestniť na príslušné miesto v texte, a prejde na novú stranu.
\end{document}